\chapter{Risk Management Frameworks}
\label{chap:Risk Management Frameworks}

%---------------------------- SECTION ----------------------------
\section{Structure Without Straitjackets}
\paragraph{If Chapter~\ref{chap:Types of Risk} gave you a map of the risk landscape, this chapter is about the different navigation systems available to help you move across it. Risk management frameworks are, at their core, structured approaches to organising and conducting risk management activity — providing a common language, a logical sequence of steps, and a set of principles to guide decisions}
\paragraph{For compliance and legal professionals, frameworks matter for several reasons. They provide \textbf{credibility} — demonstrating to regulators, auditors, and stakeholders that your organisation's approach to risk management is grounded in recognised good practice rather than invented from scratch. They provide \textbf{consistency} — ensuring that risk management activity across different teams, functions, and geographies follows a coherent and comparable approach. And they provide a \textbf{starting point} — particularly valuable for organisations that are building or maturing their risk management capability and need a reference point from which to develop.}
\paragraph{However, it is equally important to be clear about what frameworks are not. They are not prescriptive rulebooks. They are not a substitute for judgement. And adopting a framework does not, in itself, mean that your organisation is managing risk well. A framework that exists on paper but is not genuinely embedded in how the organisation operates is, at best, a missed opportunity and, at worst, a false assurance that could prove costly when things go wrong.}
\paragraph{The question this chapter seeks to answer is not \textit{which framework is correct} — there is no universally correct answer to that question. The question is: \textit{what do frameworks offer, how do the major ones differ, and how might you think about choosing an approach that genuinely serves your organisation}?}

%---------------------------- SECTION ----------------------------
\section{What a Risk Management Framework Provides}
\paragraph{Before examining specific frameworks, it is worth being clear about what a well-designed framework is intended to provide. Most recognised frameworks share a common set of aims:}
\begin{itemize}
    \bitem{A structured process}{a logical, repeatable sequence of activities for identifying, assessing, treating, and monitoring risk.}
    \bitem{Principles}{overarching guidance on the values and behaviours that should underpin risk management activity.}
    \bitem{Roles and responsibilities}{clarity about who is accountable for risk management at different levels of the organisation.}
    \bitem{Integration }{guidance on how risk management connects to governance, strategy, decision-making, and operational activity.}
    \bitem{Continuous improvement}{mechanisms for learning, reviewing, and enhancing risk management capability over time.}
    \bitem{Documentation and communication}{expectations around recording, reporting, and communicating risk information.}
\end{itemize}
\paragraph{A framework that delivers all of these things well provides a genuinely strong foundation. In practice, most organisations find that no single framework delivers everything perfectly for their specific context — and the most effective approaches tend to involve drawing on the principles of recognised frameworks whilst adapting them thoughtfully to the organisation's own circumstances.}

%---------------------------- SECTION ----------------------------
\section{ISO 31000 — The International Standard for Risk Management}
\paragraph{\textbf{ISO 31000} is the most widely recognised international standard for risk management. Published by the International Organization for Standardization and most recently revised in 2018, it provides a set of principles, a framework, and a process for managing risk that is designed to be applicable to any organisation, regardless of size, sector, or type of risk.}
\paragraph{It is important to note at the outset that \textbf{ISO 31000 is not a certifiable standard}. Unlike some other ISO standards — ISO 27001 for information security, for example — organisations cannot be formally certified against ISO 31000. It is a guidance standard, intended to inform and improve risk management practice rather than to set a compliance threshold. This distinction matters: adopting ISO 31000 is a signal of intent and approach, not a certification that can be pointed to as evidence of compliance.}
\subsection{The Principles of ISO 31000}
\paragraph{At the foundation of ISO 31000 are eight principles that describe the characteristics of effective risk management. Risk management, according to the standard, should be:}
\begin{itemize}
    \bitem{Integrated}{embedded into all organisational activities and decision-making, not treated as a standalone function.}
    \bitem{Structured and comprehensive}{following a consistent, thorough approach that produces comparable and reliable results.}
    \bitem{Customised}{tailored to the organisation's specific context, objectives, and risk profile.}
    \bitem{Inclusive }{involving appropriate stakeholders and drawing on diverse perspectives.}
    \bitem{Dynamic}{responsive to change, with the capacity to detect and adapt to new and emerging risks.}
    \bitem{Best available information}{based on the most relevant and reliable information available, whilst acknowledging that information is always imperfect.}
    \bitem{Human and cultural factors}{recognising that human behaviour and culture significantly influence risk management effectiveness.}
    \bitem{Continual improvement}{learning from experience and continuously enhancing the risk management framework and process.}
\end{itemize}
\paragraph{These principles have a distinctly non-prescriptive character — they describe \textit{what good looks like} rather than \textit{exactly how to achieve it}. This is consistent with the standard's design intent, but it also means that organisations need to do their own thinking about how these principles translate into practice in their specific context.}
\subsection{The ISO 31000 Framework}
\paragraph{The framework component of ISO 31000 describes how risk management should be organised and supported at the leadership and governance level. It emphasises leadership commitment, the integration of risk management into organisational governance, the allocation of appropriate roles and resources, and the importance of ongoing evaluation and improvement of the framework itself.}
\paragraph{For compliance professionals, the framework component is particularly relevant because it speaks directly to \textbf{governance} — who owns risk management, how it is resourced, how it connects to board-level oversight, and how the organisation ensures that its risk management capability develops over time.}
\subsection{The ISO 31000 Process}
\paragraph{The process component describes the operational steps of risk management — the activities that individuals and teams carry out to identify, assess, treat, and monitor risk. It follows a logical sequence:}
\begin{itemize}
    \bitem{Communication and consultation}{engaging stakeholders throughout the process.}
    \bitem{Scope, context, and criteria}{defining the parameters of the risk assessment.}
    \bitem{Risk assessment}{comprising risk identification, risk analysis, and risk evaluation.}
    \bitem{Risk treatment}{selecting and implementing responses to risk.}
    \bitem{Monitoring and review}{tracking the effectiveness of treatments and watching for changes in the risk landscape.}
    \bitem{Recording and reporting}{documenting the process and communicating results.}
\end{itemize}
\paragraph{We will examine each of these steps in detail in Part III. For now, the key point is that ISO 31000 provides a coherent, internationally recognised process that can serve as a useful reference point regardless of your sector or the specific nature of the risks you are managing.}

%---------------------------- SECTION ----------------------------
\section{COSO ERM — Enterprise Risk Management for Organisations}
\paragraph{\textbf{The Committee of Sponsoring Organizations of the Treadway Commission (COSO)} has developed one of the most widely adopted frameworks for enterprise risk management, particularly in North America and in organisations with a strong financial governance orientation.}
\paragraph{The most recent version of the COSO \textbf{Enterprise Risk Management — Integrating with Strategy and Performance} framework, published in 2017, makes a significant and deliberate shift from earlier versions: it places risk management firmly in the context of \textbf{strategy and performance}, rather than treating it primarily as a control and compliance activity.}
\paragraph{This is an important evolution. The 2017 COSO framework argues that the most significant risks facing organisations are often not the operational or compliance risks that traditional risk management focuses on, but the strategic risks that arise from the choices organisations make about their direction, their business model, and their performance ambitions. By integrating risk management with strategic planning and performance management, organisations are better placed to identify and manage the risks that matter most.}
\subsection{Key Features of the COSO ERM Framework}
\paragraph{The COSO ERM framework is organised around five interrelated components.}
\begin{itemize}
    \bitem{Governance and culture}{the foundation of enterprise risk management, encompassing board oversight, operating structures, and the values and behaviours that shape risk culture.}
    \bitem{Strategy and objective-setting}{integrating risk management into the strategic planning process, including the articulation of risk appetite.}
    \bitem{Performance}{identifying, assessing, and prioritising risks that affect the achievement of objectives, and implementing responses.}
    \bitem{Review and revision}{assessing the performance of risk management and adapting in response to change}
    \bitem{Information, communication, and reporting}{using risk information to support decision-making and meet reporting obligations.}
\end{itemize}
\paragraph{For compliance professionals working in organisations with a strong strategic focus — or in environments where the relationship between risk and performance is particularly visible, such as financial services — the COSO ERM framework offers a useful lens for connecting compliance activity to broader organisational objectives.}

%---------------------------- SECTION ----------------------------
\section{The Three Lines Model}
\paragraph{Formerly known as the \textbf{Three Lines of Defence}, and updated to the \textbf{Three Lines Model} by the Institute of Internal Auditors in 2020, this is less a risk management framework in its own right and more a \textbf{governance model} for organising risk management roles and responsibilities within an organisation.}
\paragraph{It is, however, one of the most widely recognised and practically useful models for compliance professionals, and it deserves careful attention.}
\subsection{The Three Lines}
\subsubsection{The First Line — Functions that own and manage risk}
\paragraph{The first line comprises the business functions and operational teams that own and manage risk as part of their day-to-day activities. These are the people closest to the risks — the teams that originate transactions, manage customer relationships, operate processes, and make operational decisions. Effective risk management in the first line means that risk awareness and risk controls are embedded in how work is done, not bolted on as an afterthought.}
\subsubsection{The Second Line — Functions that oversee risk}
\paragraph{The second line comprises the functions that provide oversight, challenge, and support to the first line — most commonly the risk management function, the compliance function, and the legal team. The second line does not own the risks themselves; its role is to ensure that the first line is identifying and managing risks appropriately, that the organisation's risk management framework is operating effectively, and that relevant standards, policies, and regulatory requirements are being met.}
\paragraph{For many readers of this book, the compliance and legal function sits squarely in the second line. This is a position of significant influence — but it comes with a corresponding responsibility to provide genuine, independent oversight and challenge, rather than simply ratifying decisions that have already been made.}
\subsubsection{The Third Line — Functions that provide independent assurance}
\paragraph{The third line is typically the internal audit function, which provides independent assurance to the board and senior management that the overall system of governance, risk management, and control is operating effectively. Unlike the second line, which is involved in designing and overseeing risk management, the third line sits entirely outside the process — its value derives from its independence}
\subsection{The ISO 31000 Process}
\paragraph{The Three Lines Model matters for compliance professionals for several reasons. It provides a clear articulation of where the compliance function sits in the governance structure, what its role is, and crucially, what it is not responsible for. One of the most common governance failures in organisations is the blurring of lines between the first and second line — where the compliance function becomes so involved in operational decision-making that it loses its independence and its ability to provide effective oversight.}
\paragraph{The model also provides a useful framework for conversations with boards and senior management about how risk management is organised and resourced — and where the gaps or weaknesses in the current structure might lie.}

%---------------------------- SECTION ----------------------------
\section{Sector-Specific Frameworks}
\paragraph{In addition to the general frameworks described above, a number of sector-specific frameworks and standards have developed that are particularly relevant in regulated industries. A selection of the most significant is outlined below}
\subsection{Basel Framework — Banking and Financial Services}
\paragraph{The \textbf{Basel Framework} — developed by the Basel Committee on Banking Supervision and currently in its third iteration (Basel III, with ongoing development towards Basel IV) — sets out international standards for bank capital adequacy, stress testing, and liquidity risk management. It provides a detailed and highly technical framework for the identification, measurement, and management of credit, market, and operational risk in banking institutions.}
\paragraph{For compliance professionals in banking or financial services, the Basel Framework is not optional reading — it underpins a significant portion of the regulatory capital and risk management requirements to which their organisations are subject. For those outside banking, it is worth understanding as a reference point, both because of its influence on regulatory thinking more broadly and because many non-bank financial institutions are subject to requirements inspired by or modelled on it.}
\subsection{FCA and PRA Frameworks — UK Financial Services}
\paragraph{In the United Kingdom, the \textbf{Financial Conduct Authority (FCA)} and the \textbf{Prudential Regulation Authority (PRA)} each publish supervisory frameworks, guidance, and expectations around risk management for the firms they regulate. The FCA's \textbf{Principles for Businesses} and \textbf{Senior Management Arrangements, Systems and Controls (SYSC)} sourcebook set out detailed expectations around risk management governance, systems, and controls.}
\paragraph{For UK-regulated firms, these are not merely guidance — they are enforceable regulatory requirements, and failure to maintain adequate risk management systems and controls is a direct basis for regulatory action}
\subsection{NIST Cybersecurity Framework}
\paragraph{The \textbf{National Institute of Standards and Technology (NIST) Cybersecurity Framework} — originally developed for US critical infrastructure but now widely adopted internationally — provides a structured approach to identifying, protecting against, detecting, responding to, and recovering from cybersecurity risks. It is increasingly referenced by regulators and organisations globally as a benchmark for cybersecurity risk management.}
\paragraph{For compliance professionals with responsibility for data protection and cybersecurity governance, the NIST framework is a valuable reference point, particularly when used alongside the risk assessment requirements of data protection legislation}
\subsection{ISO 27001 — Information Security Management}
\paragraph{\textbf{ISO 27001} is the international standard for information security management systems. Unlike ISO 31000, it is a certifiable standard — organisations can seek formal certification as evidence that their information security management meets the standard's requirements. It includes specific requirements around the identification and assessment of information security risks, making it directly relevant to the cybersecurity risk assessment work covered in Chapter~\ref{chap:Cybersecurity and Data Protection Risk Assessment}.}

%---------------------------- SECTION ----------------------------
\section{Choosing a Framework — Practical Considerations}
\paragraph{Given the range of frameworks available, how should an organisation approach the question of which to adopt? There is no universally correct answer, but the following considerations are worth working through:}
\subsection{Regulatory Expectations}
\paragraph{In some sectors and for some risk types, regulators have specific expectations about the frameworks or standards organisations should follow. Before making any other decisions, it is worth understanding whether your regulatory environment effectively mandates a particular approach — or at least strongly signals a preference for one.}
\subsection{Organisational Context and Maturity}
\paragraph{A complex, multi-jurisdictional financial institution and a mid-sized professional services firm have very different risk management needs. A framework that is appropriate and proportionate for one may be excessive or insufficient for the other. Consider your organisation's size, complexity, risk profile, and existing risk management maturity when evaluating options.}
\subsection{Alignment with Existing Governance Structures}
\paragraph{A framework that maps cleanly onto your existing governance and decision-making structures is likely to embed more effectively than one that requires significant structural change. That said, if the existing structure is itself a source of weakness, a framework adoption can be a valuable opportunity to drive improvements.}
\subsection{Practical Usability}
\paragraph{The most technically sophisticated framework is of limited value if it is too complex to be understood and applied by the people who need to use it. Practical usability — the degree to which a framework can be translated into clear, actionable guidance for the people on the ground — is an important and sometimes underweighted consideration.}
\subsection{The Option of Combining Approaches}
\paragraph{Many organisations find that no single framework meets all of their needs, and that the most effective approach involves drawing on elements of several. Using ISO 31000 as an overarching process framework, the Three Lines Model to organise governance responsibilities, and sector-specific standards for particular risk types is a perfectly reasonable approach — provided the combination is coherent and consistently applied.}


%---------------------------- SECTION ----------------------------
\section{What Frameworks Cannot Do}
\paragraph{Before closing this chapter, it is worth being explicit about the limitations of frameworks — not to discourage their use, but to ensure that expectations are realistic.}
\paragraph{A framework cannot substitute for genuine commitment to risk management at the leadership level. An organisation whose board and senior management view risk management as a compliance obligation to be discharged, rather than a genuine management tool, will not derive full value from even the best-designed framework.}
\paragraph{A framework cannot compensate for a poor risk culture. If the people working within an organisation do not feel safe raising concerns, do not believe that risk management is valued, or do not understand their own role in managing risk, the framework will remain a document rather than a reality.}
\paragraph{And a framework cannot anticipate every risk. The risks that have caused the most serious harm to organisations — and to wider society — have often been precisely those that fell outside existing frameworks, challenged existing assumptions, or emerged in ways that no one had previously considered.}
\paragraph{Frameworks are tools. Like all tools, their value depends entirely on the skill, judgement, and genuine commitment of the people who use them.}

%---------------------------- SECTION ----------------------------
\newpage
\section{Chapter Summary}
Before moving on, let us anchor the key ideas from this chapter.\\
We learned about the following regulatory frameworks:
\begin{table}[H]
  \centering
  \renewcommand{\arraystretch}{1.5}
  \begin{tabularx}{\textwidth}{ | >{\raggedright\arraybackslash}p{0.20\textwidth} 
                                | >{\raggedright\arraybackslash}p{0.30\textwidth} 
                                | >{\raggedright\arraybackslash}p{0.12\textwidth} 
                                | >{\raggedright\arraybackslash}X | }
    \hline
    \textbf{Framework / Model} & \textbf{Primary Focus} & \textbf{Certifiable?} & \textbf{Most Relevant For} \\
    \hline
    ISO31000 & General risk management principles and process & No & Any organisation seeking a recognised baseline \\
    \hline
    COSO ERM & Enterprise risk integrated with strategy and performance & No & Organisations with a strong strategic governance focus\\
    \hline
    Three Lines Model & Risk governance and accountability structure & No & Organising roles and responsibilities across the organisation\\
    \hline
    Basel Framework & Banking risk — credit, market, operational & No & Banks and regulated financial institutions\\
    \hline
    FCA / PRA Frameworks & UK financial services regulation & No & UK-regulated firms\\
    \hline
    NIST CSF & Cybersecurity risk management & No & Organisations managing cybersecurity risk\\
    \hline
    ISO 27001 & Information security management & Yes & Organisations seeking certification in information security\\
    \hline
  \end{tabularx}
  \caption{Regulatory Frameworks}
\end{table}

\paragraph{Key takeaways:}
\begin{itemize}
  \item{Frameworks provide structure, credibility, and consistency — but they are tools, not guarantees.}
  \item{ISO 31000 is the most broadly applicable international standard and a sensible reference point for most organisations.}
  \item{The Three Lines Model is particularly valuable for compliance professionals in understanding and articulating their governance role.}
  \item{Sector-specific frameworks — particularly in financial services — may effectively set the floor for regulatory expectations.}
  \item{No single framework is universally correct — the best approach is one that is appropriately tailored, genuinely embedded, and proportionate to your organisation's context.}
  \item{Frameworks cannot substitute for leadership commitment, good culture, or sound judgement}
\end{itemize}

\barquote{In Chapter~\ref{chap:Legal and Regulatory Considerations}, we move from the voluntary and guidance-based world of frameworks to the more specific and enforceable terrain of legal and regulatory requirements — exploring what compliance obligations look like across different sectors and jurisdictions, and what regulators actually expect to see when they look at your risk management arrangements.}
